\documentclass[../../../include/open-logic-section]{subfiles}

\begin{document}

\olfileid{sth}{replacement}{ref}
\olsection{પ્રતિસ્થાપન અને પ્રતિબિંબન}

\stagesinex{}ના આધારે પ્રતિસ્થાપનને વાજબી ઠરાવવાનો
આપણો છેલ્લો પ્રયાસ એક ગહન અને સુંદર પરિણામથી
શરૂ થાય છે:\footnote{યાદ રાખો: જુદું સ્પષ્ટ ન કહ્યું હોય
તો બધાં સૂત્રોમાં પરિમાણો હોઈ શકે છે.}
%In this, I will use overlining, such as $x_1, \ldots, x_n$, to
%abbreviate ``$x_1, \ldots, x_n$'': 
\begin{thm}[પ્રતિબિંબનનું પ્રરૂપ]\ollabel{reflectionschema}
કોઈ પણ સૂત્ર $\phi$ માટે:
\[
\forall \alpha \exists \beta > \alpha (\forall x_1 \ldots, x_n \in
V_\beta)(\phi(x_1, \ldots, x_n) \liff \phi^{V_\beta}(x_1, \ldots, x_n))
\]
\end{thm}
\noindent
\olref[sth][replacement][strength]{formularelativization}ની જેમ,
$\phi^{V_\beta}$ એ $\phi$ના દરેક પરિમાણકને
$V_\beta$ ગણ સુધી મર્યાદિત કરવાથી મળતું સૂત્ર છે.
તેથી સહજ રીતે, પ્રતિબિંબન કહે છે: જો $\phi$ આખા
પદાનુક્રમમાં સાચું હોય, તો $\phi$ પદાનુક્રમના ગમે તેટલા
ઊંચા \emph{આરંભિક ખંડો}માં પણ સાચું હોય છે.

\citet{Montague1961} અને \citet{Levy1960}એ બતાવ્યું
કે પ્રતિસ્થાપન અને પ્રતિબિંબનનાં (યોગ્ય રીતે ઘડેલાં)
સ્વરૂપો $\Z$ને આધાર માનીને સમકક્ષ છે. એટલે તેમાંથી
કોઈ એક ઉમેરીએ તો $\ZF$ મળે છે. (આ પરિણામો
\olref[sth][replacement][refproofs]{sec}માં સાબિત કરીશું.)
આ સમકક્ષતા હોય, તો \stagesinex{} દ્વારા પ્રતિબિંબન અને
પ્રતિસ્થાપનને આ રીતે વાજબી ઠરાવવાની આશા રાખી શકાય:
\stagesinex{} સ્વીકાર્યું હોય, તો પદાનુક્રમ ઘણો, ઘણો
ઊંચો હોવો જોઈએ; એટલો ઊંચો કે તેના વિશે આપણે કહી
શકીએ એવી કોઈ વાત તેની ઊંચાઈ મર્યાદિત કરવા પૂરતી ન હોય.
આને એવો વિચાર માની શકાય કે કોઈ પણ વિધાન~$\phi$
આખા પદાનુક્રમમાં સાચું હોય, તો તે પદાનુક્રમના ગમે
તેટલા ઊંચા આરંભિક ખંડોમાં સાચું હોય છે. અને આ જ
પ્રતિબિંબન છે.

આ ફરીથી પ્રતિસ્થાપન માટે આંતરિક ઔચિત્ય આપવાનો
ખરેખર આશાસ્પદ પ્રયાસ લાગે છે. પરંતુ તેના વિશે અહીં
કહેવાનું ઘણું બધું છે. તે સફળ થાય છે કે નહીં એ હવે
તમારે પોતે નક્કી કરવું પડશે.\footnote{જોકે આગળ
\citet[95--100]{Incurvati2020} વાંચવા જેવું લાગશે.}

\end{document}
